% Regression: p3_eq3_fusion — renders sum_m T_a^{l,m} T_b^{m,p} = sum_{c,mu} What_{ab}^{c,mu} . T_c^{l,p} .…
% Formula: sum_m T_a^{l,m} T_b^{m,p} = sum_{c,mu} What_{ab}^{c,mu} . T_c^{l,p} . W_{ab}^{c,mu}
% P3 — arXiv:2203.12563 eq (3): the fusion-tensor decomposition
%   sum_m T_a^{l,m} T_b^{m,p} = sum_{c,mu} What_{ab}^{c,mu} . T_c^{l,p} . W_{ab}^{c,mu}
% Paper ink: two red MPO wires (a over b) carrying stacked T dots with one
% physical line through both; RHS a chirality pair of fat bars (What pair
% legs west / W pair legs east) around a T_c dot on the fused wire c, with
% multiplicity mu on top of each bar.
%
% Attempt A (one picture): 3 wire rows; \tnfuse[span=3] puts the combined
%   stub exactly at row 2's height, so the T_c dot CAN sit on the fused
%   wire.  \tnskip in rows 1,3 stops the a/b wires from bridging the gap.
% Attempt B (three pictures): each bar keeps its supported orientation
%   (combined leg = the open side); the product What . T_c . W is spelled
%   as three juxtaposed math atoms sharing the row-midline axis.
\documentclass[varwidth=19cm, border=6pt]{standalone}
\usepackage{amsmath, amssymb}
\usepackage{tenkz}
\begin{document}

% --- Attempt A: single picture --------------------------------------
\[
  \begin{tenkz}[rows={op:operator, op:operator}]
    \tn{T_a} \\
    \tn{T_b}
  \end{tenkz}
  \;=\; \sum_{c,\mu}\;
  \begin{tenkz}[rows={wire:operator, op:operator, wire:operator}]
    \tnghost{} & \tnfuse[span=3, combined=east]{\mu} & \tnskip &
      \tnfuse[span=3, combined=west]{\mu} & \tnghost{} \\
               &                           & \tn{T_c} &         & \\
    \tnghost{} &                           & \tnskip  &         & \tnghost{}
  \end{tenkz}
\]

% --- Attempt B: the product as three juxtaposed atoms ----------------
\[
  \begin{tenkz}[rows={op:operator, op:operator}]
    \tn{T_a} \\
    \tn{T_b}
  \end{tenkz}
  \;=\; \sum_{c,\mu}\;
  \begin{tenkz}[rows={wire:operator, wire:operator}]
    \tnghost{} & \tnfuse[combined=east]{\mu}
  \end{tenkz}
  \begin{tenkz}[rows={op:operator}]
    \tn{T_c}
  \end{tenkz}
  \begin{tenkz}[rows={wire:operator, wire:operator}]
    \tnfuse[combined=west]{\mu} & \tnghost{}
  \end{tenkz}
\]

\end{document}
